\section{Дополнительная информация}
\label{sec:additional-information}

1. Constellation — это ряд однотипных спутников, сходных по типу и функциям, 
проектированных для размещения на сходных, взаимодополняющих орбитах ради 
общей цели, под единым управлением. Каждый спутник в составе группировки управляется 
индивидуально с Земли через наземный центр управления. Основная цель — 
обеспечение глобального или зонального покрытия для телекоммуникационных, навигационных 
или наблюдательных задач. Точность позиционирования в таких системах обычно средняя или 
низкая (порядка километров). Управление осуществляется централизованно, что позволяет 
координировать работу всех аппаратов, но требует постоянной связи с Землей. 
Классическими примерами являются системы Iridium, Globalstar и GPS, 
которые демонстрируют эффективность такого подхода для массовых коммерческих 
и государственных услуг.
Определение из источника с более 50 цитированиями с 2003 года во введении
\url{https://www.researchgate.net/publication/2559727_Satellite_Constellation_Networks}
\url{https://www.space4water.org/taxonomy/term/1679}

2. Formation — это группа спутников, в которой реализовано замкнутое управление 
на орбите через межспутниковые связи с целью сохранения заданной 
топологии и контроля относительных расстояний. Как отмечается в статье Фольта 
на странице 244 просидинга \cite{van_der_Ha_1998}, формации позволяют выполнять 
совместные измерения с высокой точностью. 
Основная задача формаций — синхронные наблюдения, интерферометрия и другие миссии, 
требующие строгого взаимного позиционирования. Управление активное, высокоточное, 
обычно с использованием бортовых систем с обратной связью. Точность позиционирования 
в формациях достигает метров, что позволяет решать задачи, недоступные для одиночных 
спутников. Примером служит миссия NASA Earth Orbiter-1 (EO-1), которая летела в 
формации с Landsat-7, повторяя его трассу с отклонением не более 3 км.

3. Swarm определяется как группа однотипных космических аппаратов, связанных между собой
для достижения общей цели без фиксированных позиций; каждый член группировки самостоятельно 
определяет и контролирует свои относительные положения относительно других. Это 
понятие тесно связано с концепцией роевого интеллекта, которая подробно рассматривается 
в работе \url{https://www.sciencedirect.com/org/science/article/pii/S1260587525000167}.
Основная цель — создание распределенных 
систем наблюдения, где большое количество недорогих спутников работает как единый 
«коллективный разум». Управление в рое децентрализованное и автономное, что 
обеспечивает высокую устойчивость к отказам и масштабируемость. Точность 
позиционирования обычно низкая и зависит от конкретной задачи, но компенсируется
избыточностью и гибкостью системы. Современные примеры — миссии HelioSwarm и 
NASA Starling, которые демонстрируют возможности роев для изучения космической 
погоды и отработки автономного управления. 
Определение содержится в абстракте \url{https://ieeexplore.ieee.org/abstract/document/5747259/citations#citations}
We define spacecraft swarms as a formation of cooperative spacecraft in which the state of each spacecraft depends on the state of the other members in the formation
из \url{https://ieeexplore.ieee.org/abstract/document/8396492/similar#similar}

4. Cluster — это распределенная гетерогенная система космических аппаратов, 
которые совместно достигают общей цели. В отличие от роя, кластер может 
включать спутники разных типов, что позволяет решать более сложные 
и разнородные задачи. Основная цель кластера — локальное усиление покрытия 
или возможностей системы, например, путем объединения спутников с различными 
сенсорами. Управление в кластере может быть смешанным, часто с выделением 
ведущего аппарата, но в рамках общей архитектуры группировки. 
Точность позиционирования в кластере — средняя, так как требования к 
взаимному расположению менее строгие, чем в формации. Примером служит 
группа спутников, одновременно находящихся в поле зрения одной наземной 
точки, что позволяет повысить пропускную способность или надежность связи.

5. Draim Orbits / Draim Constellation
Draim Orbits представляют собой отдельный класс оптимальных 
созвездий, использующих эллиптические орбиты с общим периодом и наклонением
для достижения однократного или многократного непрерывного глобального покрытия. 
Как указано на странице 53 просидинга \cite{van_der_Ha_1998}, такие созвездия 
требуют меньшего числа спутников по сравнению с круговыми орбитами. 
Основная цель — экономия аппаратов за счет вытянутых траекторий, 
которые «зависают» над нужными регионами дольше, чем круговые. 
Управление в такой группировке централизованное, с наземным контролем. 
Точность позиционирования средняя, так как эллиптические 
орбиты вносят дополнительные вариации в геометрию покрытия. 
Пример — система ELLIPSO, которая использует комбинацию эллиптических и 
круговых орбит для оптимизации покрытия северного полушария и экваториальных областей.

\begin{table}[htbp]
\centering
\caption{Сравнение терминов для групп спутников}
\label{tab:satellite_groups}
\begin{tabular}{|p{2.4cm}|p{3.0cm}|p{3.0cm}|p{2.8cm}|p{2.8cm}|p{3.0cm}|}
\hline
\textbf{Термин} & \textbf{Определение} & \textbf{Основная цель} & \textbf{Управление} & \textbf{Точность позиционирования} & \textbf{Примеры} \\
\hline
\textbf{Constellation} & Группа спутников на схожих траекториях без контроля относительного положения, каждый управляется с Земли. & Глобальное/зональное покрытие & Централизованное (наземный контроль) & Средняя/низкая (километры) & Iridium, Globalstar, GPS \\
\hline
\textbf{Formation} & Замкнутое управление через межспутниковые связи для сохранения топологии и контроля расстояний. & Совместные измерения, виртуальный инструмент & Активное, высокоточное (бортовое) & Высокая (метры) & EO-1 / Landsat-7 \\
\hline
\textbf{Swarm} & Группа однородных аппаратов, достигающих общую цель без фиксированных позиций; каждый определяет свои относительные позиции. & Распределенные системы наблюдения & Децентрализованное, автономное & Низкая (зависит от задачи) & HelioSwarm, NASA Starling \\
\hline
\textbf{Cluster} & Гетерогенная система спутников для совместного достижения общей цели. & Локальное усиление покрытия & Смешанное (лидер-ведомый) & Средняя & Спутники в поле зрения одной наземной точки \\
\hline
\textbf{Draim Constellation} & Эллиптические орбиты с общим периодом и наклонением для глобального покрытия с меньшим числом спутников. & Глобальное непрерывное покрытие & Централизованное & Средняя & ELLIPSO \\
\hline
\end{tabular}
\end{table}


Определения на английском ниже приведены из презентации 
\url{https://www.summerschoolalpbach.at/docs/2019/lectures/Schilling_Networked_Distributed_Smalll_Satellite_Systems.pdf#2#2}

constellation: similar trajectories without control for relative
position; each satellite individually controlled from ground centre.

formation: closed-loop control in orbit via inter-satellite links in
order to preserve topology in the multi-satellite system, in
particular control of relative distances

swarm: a group of similar vehicles cooperating to achieve a joint
goal without fixed positions; each member determines and
controls relative positions in relations to others.

cluster: distributed heterogeneous system of vehicles to achieve
in cooperation a joint objective
